\documentclass[10pt,a4paper]{amsart}
\usepackage{amsmath,amssymb,amsthm,graphicx}
\usepackage{verbatim}
\newenvironment{reference}{\comment}{\endcomment}
%\newenvironment{reference}{}{}
\newenvironment{slogan}{\comment}{\endcomment}
\newenvironment{history}{\comment}{\endcomment}

% For commutative diagrams we use Xy-pic
\usepackage[all]{xy}

% We use 2cell for 2-commutative diagrams.
\xyoption{2cell}
% Theorem environments.
%
\theoremstyle{plain}
\newtheorem{theorem}[subsection]{Theorem}
\newtheorem{proposition}[subsection]{Proposition}
\newtheorem{lemma}[subsection]{Lemma}

\theoremstyle{definition}
\newtheorem{definition}[subsection]{Definition}
\newtheorem{example}[subsection]{Example}
\newtheorem{exercise}[subsection]{Exercise}
\newtheorem{situation}[subsection]{Situation}

\theoremstyle{remark}
\newtheorem{remark}[subsection]{Remark}
\newtheorem{remarks}[subsection]{Remarks}

\numberwithin{equation}{subsection}

% Macros
%
\def\lim{\mathop{\mathrm{lim}}\nolimits}
\def\colim{\mathop{\mathrm{colim}}\nolimits}
\def\Spec{\mathop{\mathrm{Spec}}}
\def\Hom{\mathop{\mathrm{Hom}}\nolimits}
\def\Ext{\mathop{\mathrm{Ext}}\nolimits}
\def\SheafHom{\mathop{\mathcal{H}\!\mathit{om}}\nolimits}
\def\SheafExt{\mathop{\mathcal{E}\!\mathit{xt}}\nolimits}
\def\Sch{\mathit{Sch}}
\def\Mor{\mathop{\mathrm{Mor}}\nolimits}
\def\Ob{\mathop{\mathrm{Ob}}\nolimits}
\def\Sh{\mathop{\mathit{Sh}}\nolimits}
\def\NL{\mathop{N\!L}\nolimits}
\def\CH{\mathop{\mathrm{CH}}\nolimits}
\def\proetale{{pro\text{-}\acute{e}tale}}
\def\etale{{\acute{e}tale}}
\def\QCoh{\mathit{QCoh}}
\def\Ker{\mathop{\mathrm{Ker}}}
\def\Im{\mathop{\mathrm{Im}}}
\def\Coker{\mathop{\mathrm{Coker}}}
\def\Coim{\mathop{\mathrm{Coim}}}

% Boxtimes
%
\DeclareMathSymbol{\boxtimes}{\mathbin}{AMSa}{"02}

%
% Macros for moduli stacks/spaces
%
\def\QCohstack{\mathcal{QC}\!\mathit{oh}}
\def\Cohstack{\mathcal{C}\!\mathit{oh}}
\def\Spacesstack{\mathcal{S}\!\mathit{paces}}
\def\Quotfunctor{\mathrm{Quot}}
\def\Hilbfunctor{\mathrm{Hilb}}
\def\Curvesstack{\mathcal{C}\!\mathit{urves}}
\def\Polarizedstack{\mathcal{P}\!\mathit{olarized}}
\def\Complexesstack{\mathcal{C}\!\mathit{omplexes}}
% \Pic is the operator that assigns to X its picard group, usage \Pic(X)
% \Picardstack_{X/B} denotes the Picard stack of X over B
% \Picardfunctor_{X/B} denotes the Picard functor of X over B
\def\Pic{\mathop{\mathrm{Pic}}\nolimits}
\def\Picardstack{\mathcal{P}\!\mathit{ic}}
\def\Picardfunctor{\mathrm{Pic}}
\def\Deformationcategory{\mathcal{D}\!\mathit{ef}}
\usepackage{cleveref}
\crefname{originallemma}{original lemma}{original lemmas}
\Crefname{originallemma}{Original lemma}{Original lemmas}
\makeatletter
\newlabel{algebra-lemma-completion-flat}{{00MB}{source}{}{sourcepointer-1}{}}
\newlabel{algebra-lemma-completion-flat@cref}{{[originallemma][0][]00MB}{[1][0][]source}{}{}{}}
\makeatother
\crefname{originallemma}{original lemma}{original lemmas}
\Crefname{originallemma}{Original lemma}{Original lemmas}
\makeatletter
\newlabel{coherent-lemma-flat-base-change-cohomology}{{02KH}{source}{}{sourcepointer-2}{}}
\newlabel{coherent-lemma-flat-base-change-cohomology@cref}{{[originallemma][0][]02KH}{[1][0][]source}{}{}{}}
\makeatother
\crefname{originalreference}{original reference}{original references}
\Crefname{originalreference}{Original reference}{Original references}
\makeatletter
\newlabel{lemma-blow-up-normal-rational}{{0B4Z}{source}{}{sourcepointer-3}{}}
\newlabel{lemma-blow-up-normal-rational@cref}{{[originalreference][0][]0B4Z}{[1][0][]source}{}{}{}}
\makeatother
\crefname{originalreference}{original reference}{original references}
\Crefname{originalreference}{Original reference}{Original references}
\makeatletter
\newlabel{lemma-descend-admissible-blowup}{{0AFK}{source}{}{sourcepointer-4}{}}
\newlabel{lemma-descend-admissible-blowup@cref}{{[originalreference][0][]0AFK}{[1][0][]source}{}{}{}}
\makeatother
\crefname{originalreference}{original reference}{original references}
\Crefname{originalreference}{Original reference}{Original references}
\makeatletter
\newlabel{lemma-dominate-by-scheme-modification}{{0B4M}{source}{}{sourcepointer-5}{}}
\newlabel{lemma-dominate-by-scheme-modification@cref}{{[originalreference][0][]0B4M}{[1][0][]source}{}{}{}}
\makeatother
\crefname{originalreference}{original reference}{original references}
\Crefname{originalreference}{Original reference}{Original references}
\makeatletter
\newlabel{lemma-exact-sequence}{{0AXF}{source}{}{sourcepointer-6}{}}
\newlabel{lemma-exact-sequence@cref}{{[originalreference][0][]0AXF}{[1][0][]source}{}{}{}}
\makeatother
\crefname{originalreference}{original reference}{original references}
\Crefname{originalreference}{Original reference}{Original references}
\makeatletter
\newlabel{lemma-iso-completions}{{0BG5}{source}{}{sourcepointer-7}{}}
\newlabel{lemma-iso-completions@cref}{{[originalreference][0][]0BG5}{[1][0][]source}{}{}{}}
\makeatother
\crefname{originallemma}{original lemma}{original lemmas}
\Crefname{originallemma}{Original lemma}{Original lemmas}
\makeatletter
\newlabel{more-algebra-lemma-finite-type-over-excellent}{{07QU}{source}{}{sourcepointer-8}{}}
\newlabel{more-algebra-lemma-finite-type-over-excellent@cref}{{[originallemma][0][]07QU}{[1][0][]source}{}{}{}}
\makeatother
\crefname{originallemma}{original lemma}{original lemmas}
\Crefname{originallemma}{Original lemma}{Original lemmas}
\makeatletter
\newlabel{more-algebra-lemma-normal-goes-up}{{0BFK}{source}{}{sourcepointer-9}{}}
\newlabel{more-algebra-lemma-normal-goes-up@cref}{{[originallemma][0][]0BFK}{[1][0][]source}{}{}{}}
\makeatother
\crefname{originalproposition}{original proposition}{original propositions}
\Crefname{originalproposition}{Original proposition}{Original propositions}
\makeatletter
\newlabel{more-algebra-proposition-ubiquity-excellent}{{07QW}{source}{}{sourcepointer-10}{}}
\newlabel{more-algebra-proposition-ubiquity-excellent@cref}{{[originalproposition][0][]07QW}{[1][0][]source}{}{}{}}
\makeatother
\crefname{originallemma}{original lemma}{original lemmas}
\Crefname{originallemma}{Original lemma}{Original lemmas}
\makeatletter
\newlabel{varieties-lemma-modification-normal-iso-over-codimension-1}{{0BFP}{source}{}{sourcepointer-11}{}}
\newlabel{varieties-lemma-modification-normal-iso-over-codimension-1@cref}{{[originallemma][0][]0BFP}{[1][0][]source}{}{}{}}
\makeatother
\title{Rationality of a normal surface singularity passes to the completion, and the blowup of the maximal ideal has normal completions}
\author{The Stacks Project Authors}
\date{stacks-project:resolve.tex:local-snapshot-20261012; GNU FDL 1.2 or later}
\begin{document}
\maketitle
\noindent\textit{Editorial note.} All delivered mathematical text is the author's own, copied from the pinned Stacks Project chapter file resolve.tex; no statement or proof step was rewritten, reordered or removed.\par\medskip
\noindent\textit{Editorial note.} The counted proof is the author's complete argument: the reduction of the first assertion to a vanishing statement with the spreading of admissible blowups and the normalization step, the transfer of the vanishing along the completion by flat base change, and the second assertion through normality of the blowups, excellency and the isomorphism of the relevant completions.\par\medskip
\noindent\textit{Editorial note.} The label lemma-blowup-still-good names the counted statement, and the proof environment immediately following it in the source is delivered verbatim as its proof.\par\medskip
\noindent\textit{Editorial note.} The counted statement is carried by the author's own lemma environment in the chapter source.\par\medskip
\noindent\textit{Editorial note.} Pre-open gate: python operations/agent-authoring/tools/precard\_gate.py --source resolve.tex --label lemma-blowup-still-good --start 3052 --end 3110 --local-id 28079 returned RESULT=OK (L28079 free, occupancy FREE, public interval clear). Pre-assembly tag self-check: canonical tag resolved and not present in reference/public-items.jsonl.\par\medskip
\noindent\textit{Editorial note.} The statement and proof are copied from the pinned chapter file; the cited inputs are kept as external pointers with their canonical Stacks tags.\par\medskip
\noindent\textit{Allowed standard prerequisites (editorial).}\begin{itemize}
\item Isomorphism locus of a modification of normal surfaces. Varieties, the lemma that a modification of a normal surface is an isomorphism over the complement of a finite set of closed points. Supplies the identification of the modification with an isomorphism over the punctured spectrum of the completion.
\item Spreading of admissible blowups and their descent. Same chapter, the lemnas producing an admissible blowup dominating a given modification over the completion and descending admissible blowups from the completion to the original ring. Supply the admissible blowup over the original ring whose base change dominates the given modification.
\item Flat base change for the cohomology of the structure sheaf and normality of the blowup. Cohomology of Schemes, the flat base change lemma together with the flatness of completion; Same chapter, the lemma that the blowup of the maximal ideal of a normal surface with rational singularity is normal. Supply the vanishing of the first cohomology after base change and the normality of the blowups involved in the second assertion.
\item Excellency of complete local rings and regularity of the completion map. More on Algebra, the ubiquity of excellent rings, the lemma that a finite type algebra over an excellent ring is excellent, and the lemma that normality ascends along a regular homomorphism. Supply the normality of the completed local rings of the blowup.
\end{itemize}\medskip
\ifdefined\hypertarget\hypertarget{sourcepointer-1}{}\fi
\noindent\textit{Source reference.} Original lemma 00MB: sources/stacks-chow/tags.txt, source lines 817--817. Prerequisite pointer invoked at proof line(s) 3085; it names a standard fact of a sibling chapter that the counted argument uses, and every step of the counted proof built on it is retained verbatim in the delivered body.\par\medskip
\ifdefined\hypertarget\hypertarget{sourcepointer-2}{}\fi
\noindent\textit{Source reference.} Original lemma 02KH: sources/stacks-chow/tags.txt, source lines 3372--3372. Prerequisite pointer invoked at proof line(s) 3083; it names a standard fact of a sibling chapter that the counted argument uses, and every step of the counted proof built on it is retained verbatim in the delivered body.\par\medskip
\ifdefined\hypertarget\hypertarget{sourcepointer-3}{}\fi
\noindent\textit{Source reference.} Original reference 0B4Z: sources/stacks-resolve/pinned-original-source.tex, source lines 2534--2534. Prerequisite pointer invoked at proof line(s) 3093; it names a standard fact of the same chapter that the counted argument uses, and every step of the counted proof built on it is retained verbatim in the delivered body.\par\medskip
\ifdefined\hypertarget\hypertarget{sourcepointer-4}{}\fi
\noindent\textit{Source reference.} Original reference 0AFK: sources/stacks-resolve/pinned-original-source.tex, source lines 3029--3029. Prerequisite pointer invoked at proof line(s) 3072; it names a standard fact of the same chapter that the counted argument uses, and every step of the counted proof built on it is retained verbatim in the delivered body.\par\medskip
\ifdefined\hypertarget\hypertarget{sourcepointer-5}{}\fi
\noindent\textit{Source reference.} Original reference 0B4M: sources/stacks-resolve/pinned-original-source.tex, source lines 1612--1612. Prerequisite pointer invoked at proof line(s) 3070; it names a standard fact of the same chapter that the counted argument uses, and every step of the counted proof built on it is retained verbatim in the delivered body.\par\medskip
\ifdefined\hypertarget\hypertarget{sourcepointer-6}{}\fi
\noindent\textit{Source reference.} Original reference 0AXF: sources/stacks-resolve/pinned-original-source.tex, source lines 1901--1901. Prerequisite pointer invoked at proof line(s) 3087; it names a standard fact of the same chapter that the counted argument uses, and every step of the counted proof built on it is retained verbatim in the delivered body.\par\medskip
\ifdefined\hypertarget\hypertarget{sourcepointer-7}{}\fi
\noindent\textit{Source reference.} Original reference 0BG5: sources/stacks-resolve/pinned-original-source.tex, source lines 2966--2966. Prerequisite pointer invoked at proof line(s) 3108; it names a standard fact of the same chapter that the counted argument uses, and every step of the counted proof built on it is retained verbatim in the delivered body.\par\medskip
\ifdefined\hypertarget\hypertarget{sourcepointer-8}{}\fi
\noindent\textit{Source reference.} Original lemma 07QU: sources/stacks-chow/tags.txt, source lines 9844--9844. Prerequisite pointer invoked at proof line(s) 3099; it names a standard fact of a sibling chapter that the counted argument uses, and every step of the counted proof built on it is retained verbatim in the delivered body.\par\medskip
\ifdefined\hypertarget\hypertarget{sourcepointer-9}{}\fi
\noindent\textit{Source reference.} Original lemma 0BFK: sources/stacks-chow/tags.txt, source lines 14388--14388. Prerequisite pointer invoked at proof line(s) 3103; it names a standard fact of a sibling chapter that the counted argument uses, and every step of the counted proof built on it is retained verbatim in the delivered body.\par\medskip
\ifdefined\hypertarget\hypertarget{sourcepointer-10}{}\fi
\noindent\textit{Source reference.} Original proposition 07QW: sources/stacks-chow/tags.txt, source lines 9846--9846. Prerequisite pointer invoked at proof line(s) 3096; it names a standard fact of a sibling chapter that the counted argument uses, and every step of the counted proof built on it is retained verbatim in the delivered body.\par\medskip
\ifdefined\hypertarget\hypertarget{sourcepointer-11}{}\fi
\noindent\textit{Source reference.} Original lemma 0BFP: sources/stacks-chow/tags.txt, source lines 14392--14392. Prerequisite pointer invoked at proof line(s) 3067; it names a standard fact of a sibling chapter that the counted argument uses, and every step of the counted proof built on it is retained verbatim in the delivered body.\par\medskip
\begin{lemma}
\label{lemma-blowup-still-good}
Let $(A, \mathfrak m, \kappa)$ be a Nagata local normal domain of
dimension $2$. Assume $A$ defines a rational singularity and that
the completion $A^\wedge$ of $A$ is normal. Then
\begin{enumerate}
\item $A^\wedge$ defines a rational singularity, and
\item if $X \to \Spec(A)$ is the blowing up in $\mathfrak m$, then
for a closed point $x \in X$ the completion $\mathcal{O}_{X, x}$ is normal.
\end{enumerate}
\end{lemma}

\begin{proof}
Let $Y \to \Spec(A^\wedge)$ be a modification with $Y$ normal.
We have to show that $H^1(Y, \mathcal{O}_Y) = 0$. By Varieties, Lemma
\ref{varieties-lemma-modification-normal-iso-over-codimension-1}
$Y \to \Spec(A^\wedge)$ is an isomorphism over the punctured
spectrum $U^\wedge = \Spec(A^\wedge) \setminus \{\mathfrak m^\wedge\}$.
By Lemma \ref{lemma-dominate-by-scheme-modification}
there exists a $U^\wedge$-admissible blowup $Y' \to \Spec(A^\wedge)$
dominating $Y$. By Lemma \ref{lemma-descend-admissible-blowup}
we find there exists a $U$-admissible blowup $X \to \Spec(A)$
whose base change to $A^\wedge$ dominates $Y$.
Since $A$ is Nagata, we can replace $X$ by its normalization
after which $X \to \Spec(A)$ is a normal modification (but
possibly no longer a $U$-admissible blowup).
Then $H^1(X, \mathcal{O}_X) = 0$ as $A$ defines a rational
singularity. It follows that
$H^1(X \times_{\Spec(A)} \Spec(A^\wedge),
\mathcal{O}_{X \times_{\Spec(A)} \Spec(A^\wedge)}) = 0$
by flat base change (Cohomology of Schemes, Lemma
\ref{coherent-lemma-flat-base-change-cohomology}
and flatness of $A \to A^\wedge$ by
Algebra, Lemma \ref{algebra-lemma-completion-flat}).
We find that $H^1(Y, \mathcal{O}_Y) = 0$ by
Lemma \ref{lemma-exact-sequence}.

\medskip\noindent
Finally, let $X \to \Spec(A)$ be the blowing up of $\Spec(A)$
in $\mathfrak m$. Then $Y = X \times_{\Spec(A)} \Spec(A^\wedge)$
is the blowing up of $\Spec(A^\wedge)$ in $\mathfrak m^\wedge$.
By Lemma \ref{lemma-blow-up-normal-rational} we see that both $Y$
and $X$ are normal. On the other hand, $A^\wedge$ is excellent
(More on Algebra, Proposition
\ref{more-algebra-proposition-ubiquity-excellent})
hence every affine open in $Y$ is the spectrum of an
excellent normal domain
(More on Algebra, Lemma \ref{more-algebra-lemma-finite-type-over-excellent}).
Thus for $y \in Y$ the ring map
$\mathcal{O}_{Y, y} \to \mathcal{O}_{Y, y}^\wedge$
is regular and by
More on Algebra, Lemma \ref{more-algebra-lemma-normal-goes-up}
we find that $\mathcal{O}_{Y, y}^\wedge$ is normal.
If $x \in X$ is a closed point of the special fibre,
then there is a unique closed point $y \in Y$ lying over $x$.
Since $\mathcal{O}_{X, x} \to \mathcal{O}_{Y, y}$ induces
an isomorphism on completions (Lemma \ref{lemma-iso-completions})
we conclude.
\end{proof}

\end{document}
